\section{Capítulo~\ref{sec:acet}. Acrescentamos o \nt{ACET}}

As três ações da acetilcolina foram medidas em fatias e no cérebro vivo, o conflito entre escrita e leitura foi mostrado no modelo do livro, e que o \nt{ACET} sirva de linha de habilitação de escrita e informe a incerteza esperada é uma hipótese com apoio experimental parcial.

{\footnotesize
\setlength{\tabcolsep}{3pt}\renewcommand{\arraystretch}{1.15}
\begin{longtable}{>{\raggedright}p{0.5cm}>{\raggedright}p{4.0cm}>{\raggedright}p{5.6cm}>{\raggedright}p{2.3cm}>{\raggedright\arraybackslash}p{1.7cm}}
\toprule
N.º & Proposição & Experimento e o que foi medido & Fonte & Status \\
\midrule\endfirsthead
\toprule
N.º & Proposição & Experimento e o que foi medido & Fonte & Status \\
\midrule\endhead
1 & Os neurônios \nt{ACET} do cérebro são poucos, estão reunidos em alguns grupos, e as saídas se espalham por todo o córtex: um canal de difusão & Coloração com anticorpos contra a enzima de síntese da acetilcolina e marcação retrógrada no rato: córtex, hipocampo, amígdala, bulbo olfativo e tálamo recebem acetilcolina principalmente de seis grupos compactos de células do prosencéfalo basal e do tronco & \cite{Mesulam1983} & medido \\
2 & O nível de acetilcolina no córtex é alto na vigília e baixo no sono de ondas lentas & Microdiálise no córtex e no hipocampo de gatos em livre movimento, com registro de EEG: a liberação de acetilcolina na vigília tranquila é $\sim100\%$ e na ativa $\sim175\%$ maior que no sono de ondas lentas; no sono REM, no córtex, é como na vigília & \cite{Marrosu1995,Hasselmo1999} & medido \\
3 & O sentido do sinal é definido pelo receptor: a acetilcolina tem receptores-canais rápidos e receptores lentos, que disparam reações na célula (proposição~\ref{prop:receptor}) & Revisão de medições: o efeito da acetilcolina sobre excitabilidade, transmissão e plasticidade depende do local de liberação, do subtipo de receptor (nicotínicos: canais iônicos; muscarínicos: via proteínas G) e da célula destinatária & \cite{Picciotto2012} & medido \\
4 & Primeira ação: enfraquecer as conexões internas, quase sem afetar as entradas externas (fator $1-a$ em~\eqref{eq:acetnet}) & Fatias de córtex olfativo de rato: com $100$\,\textmu M de agonistas da acetilcolina, os potenciais das conexões internas (camada Ib) caem mais de $60\%$, e os da entrada externa (camada Ia), menos de $15\%$, na mesma célula; a supressão é pré-sináptica. Fatias de hipocampo (CA1): $89\%$ contra $40\%$ para as vias interna e de entrada & \cite{HasselmoBower1992,HasselmoSchnell1994} & medido \\
5 & Segunda ação: reforçar a entrada (fator $\frac{1+a}{2}$) & Fatias de neocórtex: receptores nicotínicos reforçam as sinapses da entrada vinda do tálamo e não agem sobre as intracorticais; os muscarínicos enfraquecem ambos os tipos. A acetilcolina aumenta a excitabilidade dos neurônios (revisão) & \cite{Gil1997,Hasselmo2006} & medido \\
6 & Terceira ação: facilitar a mudança dos pesos (taxa $\eta a$) & Rato: a estimulação elétrica do núcleo basal (fonte de acetilcolina para o córtex) pareada com um tom provoca, no animal adulto, uma forte reorganização do mapa do córtex auditivo em favor do som apresentado & \cite{KilgardMerzenich1998,Hasselmo2006} & medido \\
7 & A regra de Hebb para padrões esparsos na segunda equação de~\eqref{eq:acetnet} dá uma memória associativa que completa o padrão a partir de uma parte & Cálculo da capacidade de uma rede com padrões esparsos e regra de covariância: com fração pequena $p$ de neurônios ativos, a rede guarda muito mais padrões que com $p=1/2$ & \cite{Tsodyks1988} & teoria \\
8 & Escrever com $a$ baixo corrompe a memória: a rede grava o lembrado, e não o visto; com $a=0.1$, a corrupção não é menor que com $a=0$ & \texttt{models/\allowbreak acet\_memory.py}: $200$ neurônios, cinco padrões de $20$, um padrão novo compartilha $10$ neurônios com um deles, $10$ execuções. Com $a=0$, o padrão novo não é memorizado, e o antigo arrasta $5.9$ neurônios alheios de $10$; com $a=0.1$, o novo é lembrado ($0.97$), mas os dois se fundem: o novo arrasta $7.3$ neurônios alheios, o antigo $7.9$; a partir de $a=0.2$, menos de um & dedução no capítulo & modelo \\
9 & Proposição~\ref{prop:writeread}: não há nível $a$ em que escrita e leitura sejam igualmente boas (fig.~\ref{fig:acetsim}) & Mesmo script, taxa de aprendizado $\eta a$: o padrão novo é memorizado por inteiro numa apresentação com $a\ge0.9$, em $0.8$ com $a=0.75$, e não é memorizado com $a\le0.5$. Leitura pela metade do padrão: $1.0$ com $a\le0.2$, $0.96$ com $a=0.3$, $0.04$ com $a=0.4$ & dedução no capítulo & modelo \\
10 & No cérebro, um único fio de difusão, o \nt{ACET}, comuta escrita e leitura & A ideia vem de modelos construídos a partir de medições em fatias. O experimento mais direto é em humanos: o bloqueador de receptores muscarínicos escopolamina prejudica a memorização de novos pares de palavras, sobretudo os que se sobrepõem aos antigos, mas não a lembrança de pares já aprendidos. Não há medição direta dos dois modos da rede & \cite{HasselmoAndersonBower1992,Atri2004} & hipótese \\
11 & O filtro~\eqref{eq:kalman} é ótimo; o peso da entrada e a taxa de aprendizado são a mesma grandeza $P/R$; em regime estacionário, $P=\sqrt{qR}$ e a memória é $\sqrt{R/q}$ (proposição~\ref{prop:kalman}) & Não requer experimento: a otimalidade do filtro de Kalman--Bucy é um resultado clássico da teoria de estimação; o regime estacionário é deduzido no capítulo & \cite{KalmanBucy1961} & teoria \\
12 & O \nt{ACET} informa à rede uma grandeza semelhante a $P$, a incerteza esperada & Proposto num modelo probabilístico de atenção. Não há medição direta de que o nível de acetilcolina acompanhe a incerteza da estimativa & \cite{YuDayan2005} & hipótese \\
13 & Parte dos neurônios \nt{ACET} responde à recompensa e à punição em duas dezenas de milissegundos, tanto mais forte quanto mais inesperado o reforço & Camundongo, neurônios colinérgicos do prosencéfalo basal identificados optogeneticamente: resposta à recompensa e à punição após $18\pm3$\,ms, cuja magnitude cresce com a surpresa do reforço; as respostas são semelhantes nos neurônios de dois núcleos & \cite{Hangya2015} & medido \\
14 & O \nt{NORA} percebe uma mudança inesperada do mundo; o \nt{ACET} mantém a rede no modo de escrita & A divisão de trabalho foi proposta num modelo. Indiretamente: em humanos, o diâmetro da pupila, ligado ao \nt{NORA}, acompanha as mudanças e a taxa de aprendizado. Não há teste direto do par \nt{NORA}--\nt{ACET} & \cite{YuDayan2005,Nassar2012} & hipótese \\
15 & No sono de ondas lentas, a rede em modo de leitura reproduz o que foi gravado de dia, e essas repetições o transcrevem para a memória de longo prazo & Registros de conjuntos de neurônios do hipocampo de rato: células que funcionaram juntas de dia disparam juntas com mais frequência no sono de ondas lentas seguinte, e na mesma ordem: medido. Quanto à transcrição: suprimir as ondulações (ripples) do hipocampo após o aprendizado prejudica a memória espacial, e prolongar as ondulações espontâneas a melhora; que a causa seja o nível baixo de \nt{ACET} não foi demonstrado & \cite{WilsonMcNaughton1994,Skaggs1996,Girardeau2009,FernandezRuiz2019,Hasselmo1999} & hipótese \\
\bottomrule
\end{longtable}}
